\documentclass[11pt,letterpaper]{article}
\usepackage[T1]{fontenc}
\usepackage[utf8]{inputenc}
\usepackage{lmodern,microtype,amsmath,booktabs,tabularx,longtable,graphicx}
\usepackage[margin=0.8in,headheight=15pt]{geometry}
\usepackage{xcolor,fancyhdr,enumitem,xurl}
\usepackage[hidelinks]{hyperref}
\definecolor{Gold}{HTML}{987100}
\setlength{\parindent}{0pt}\setlength{\parskip}{6pt}
\setlist{nosep,leftmargin=*}\emergencystretch=3em
\pagestyle{fancy}\fancyhf{}\fancyhead[L]{\small Follow My Voice}\fancyfoot[R]{\thepage}
\newcommand{\qid}[1]{\par\smallskip\textbf{\texttt{#1}}\par}
\begin{document}
\setcounter{section}{4}
\section{Planned Analysis}
\subsection{Primary contrast}
The primary estimand will compare active completion time under the self-similar and participant-preferred catalog voices, standardized to the same starting-difficulty distribution. A trial-level mixed-effects analysis of log completion time will account for participant clustering, voice condition, period/order, trial progression and absolute pre-cue starting theta. Report the voice time ratio and uncertainty using a two-sided contrast.

A prespecified low-complexity function of theta will allow nonlinearity; nominal-angle strata will support balance plots and a categorical sensitivity analysis. Zero-degree trials will remain in the analysis, and completion time will not be divided by angle. The model's random-effects structure, angle function and smallest meaningful effect will be finalized with simulation before confirmatory recruitment. Participant count will be justified under that same model; repeated trials will not be treated as independent participants.

\subsection{Secondary and exploratory outcomes}
Retrospective condition-specific ratings will be collected only after both blocks, with actual survey timestamps retained. Reference scope remains participant-by-block. Recall and recency effects will be considered when interpreting condition differences; these are not immediate post-block measurements. Individual workload dimensions and complete unweighted workload means will be summarized by participant and block. IMI subscales and the three six-item agent-perception constructs will remain separate scores. The additional Guo Q20--22 appealing, eerie and familiar ratings will be analyzed individually, without a combined eeriness score. Voice similarity will be a manipulation check. Custom experience ratings and self-reported guidance use will be analyzed individually. Multiplicity families and secondary contrasts will be specified before confirmatory analysis.

Instruction count, observed gaze path and other behavior potentially caused by voice condition will not automatically be covariates in the primary model. Voice-by-theta variation, selected catalog voice, acoustic quality and guidance timing will be exploratory or sensitivity analyses with their assumptions reported.

\subsection{Incomplete and interrupted observations}
Raw attempts, incorrect selections, pause intervals, timeouts, technical failures and withdrawals will be retained. Stopped trials will not be coded as successful completion times. Timeout/censoring rules and the missing-data estimand will be specified before collection; slow trials will not be removed solely for being slow. Missing responses will not be replaced with zero. Questionnaire scores will follow the complete-subscale rule, with item-level missingness reported. Disliking or failing to recognize a cloned voice will not by itself exclude a participant.

\subsection{Guidance and qualitative evaluation}
Engineering analysis will distinguish event qualification, decision latency, model deferral, scheduling delay and actual playback. Gaze and controller confirmation will be evaluated separately. Sequence-level review will include what happened before and after a cue. Automatic checks and model-judge ratings will remain advisory engineering evidence, not participant perception scores.

Debrief responses will be coded with a documented thematic process and an audit trail, with independent review of a subset. Perceived performance, voice comfort, reliance and possible own-voice use cases will be interpreted alongside, but not substituted for, measured task performance.
\end{document}
